\documentclass[11pt]{article}

\usepackage[a4paper,margin=1in]{geometry}
\usepackage{microtype}
\usepackage{amsmath}
\usepackage{booktabs}
\usepackage{enumitem}
\usepackage{titlesec}
\usepackage{xcolor}
\usepackage[hidelinks]{hyperref}
\usepackage{parskip}

\definecolor{accent}{HTML}{3C6B4C}
\titleformat{\section}{\large\bfseries\color{accent}}{\thesection}{0.6em}{}
\titleformat{\subsection}{\normalsize\bfseries}{\thesubsection}{0.6em}{}
\hypersetup{colorlinks=true, linkcolor=accent, urlcolor=accent, citecolor=accent}

\newcommand{\eps}{\varepsilon}
\newcommand{\code}[1]{\texttt{#1}}

\title{\vspace{-2em}\textbf{Attested Test Selection with a Verifiable Confidence Bound}\\
\large A validation-intelligence layer for continuous integration}
\author{Amit Patole\\
\small \href{https://orcid.org/0009-0009-6514-0913}{ORCID: 0009-0009-6514-0913}\\
\small \texttt{amit.patole@gmail.com}}
\date{July 2026}

\begin{document}
\maketitle

\begin{abstract}
\noindent
Regression test suites grow until running them on every change becomes the
dominant cost of continuous integration (CI). Test-selection and
test-prioritization tools reduce that cost by running a subset of the suite, but
they return a \emph{speedup} and ask the engineer to trust that the omitted
tests would not have caught anything. We present \textbf{assaylab}, a
validation-intelligence layer that reduces a suite \emph{and} emits a signed,
independently re-derivable receipt that \emph{bounds the confidence lost}: given
a subset $S$ and a skipped set $U$, the receipt commits to
$\eps = 1 - \prod_{t \in U}(1 - q_t)$, an upper bound on the probability that at
least one skipped test would have failed, where $q_t$ is each test's modelled
detection probability estimated from history. The bound is bound into an
HMAC-SHA256 or ed25519 signature over the actual outcome, so a consumer can
verify authenticity and, because selection is deterministic in its committed
inputs, recompute the bound to confirm it is genuine rather than merely signed.
assaylab additionally clusters failures into root-cause signatures and
distinguishes flaky failures from real ones, feeding the per-test probabilities
the bound relies on. On the FlakeFlagger public dataset (26{,}765 tests, each
rerun 10{,}000 times), using the \emph{measured} per-run failure rates as $q_t$,
assaylab selected a subset achieving a \textbf{35.9$\times$} reduction in
test-time while the realized regression-miss rate ($0.049957$) stayed within the
bound the receipt claimed ($0.049958$). The implementation is open source, was
hardened over five independent adversarial review rounds, and is available on
PyPI.
\end{abstract}

\section{Introduction}

As a codebase and its test suite grow, running every test on every change
becomes the bottleneck of CI. The standard responses --- regression test
selection (RTS) and test-case prioritization (TCP) --- run only a subset of the
suite, trading coverage for speed. The engineering objection is trust: when a
tool skips 90\% of the suite and the build is green, what is the actual
probability that a regression slipped through the skipped tests? Existing tools
report a speedup and a heuristic; they do not attach a checkable, quantitative
statement of the risk incurred, and they certainly do not attach a
\emph{verifiable} one that a downstream consumer --- a release gate, an auditor,
a compliance process --- can confirm without re-running the analysis or trusting
the producer.

This paper describes \textbf{assaylab}, whose distinctive contribution is
\emph{attested test selection with a verifiable confidence bound}. Concretely:

\begin{enumerate}[leftmargin=1.4em,itemsep=2pt]
  \item A selection engine that chooses a subset under a target confidence loss
        $\eps$ or a time budget and reports the achieved $\eps$
        (Section~\ref{sec:selection}).
  \item A signed receipt that binds the \emph{actual outcome} --- the committed
        inputs, the selected and skipped sets, and $\eps$ --- so the signature
        certifies what happened, not merely that a run happened
        (Section~\ref{sec:attest}).
  \item Independent re-derivation: because selection is deterministic in its
        committed inputs, a verifier recomputes $\eps$ and confirms the receipt
        reproduces (Section~\ref{sec:attest}).
  \item An evidence layer --- failure-signature clustering, root-cause
        categorization, and flaky-vs-real classification --- that supplies the
        per-test detection probabilities the bound consumes
        (Section~\ref{sec:evidence}).
\end{enumerate}

We evaluate the bound on real data (Section~\ref{sec:eval}) and report an
adversarial security review of the signing and gating machinery
(Section~\ref{sec:security}), then state the assumptions the bound rests on
honestly (Section~\ref{sec:limits}).

\section{Background and Related Work}

\textbf{Flaky tests.} A flaky test passes and fails without any code change.
iDFlakies~\cite{idflakies} framed detection-by-rerunning and produced the IDoFT
category registry; FlakeFlagger~\cite{flakeflagger} predicts flakiness from
static and dynamic features \emph{without} rerunning, and ships a dataset in
which each of 26{,}765 tests was rerun 10{,}000 times to establish ground truth.
DeFlaker~\cite{deflaker} detects flakiness via change-coverage. We use the
FlakeFlagger dataset both for flaky prediction and, novelly, as a source of
measured per-test failure rates against which to validate our confidence bound.

\textbf{Regression test selection / prioritization.} RTS and TCP have a large
literature; RTPTorrent~\cite{rtptorrent} provides commit-linked test-execution
histories for evaluating prioritization, and TravisTorrent~\cite{travistorrent}
supplies broad CI build outcomes. These works optimize \emph{which} tests to run
and \emph{in what order}; to our knowledge none attach a signed, re-derivable
bound on the confidence lost by the tests they skip. That combination ---
selection plus an attested, checkable risk statement --- is the gap assaylab
addresses.

\section{The Evidence Layer}
\label{sec:evidence}

assaylab ingests test outcomes (JUnit/xUnit XML, or an outcome
CSV/JSON schema of \code{(project, commit, test\_id, run\_id, verdict,
duration)}) and produces a graded \emph{report} on a shared verdict contract
(pass/warn/fail plus grounded issues).

\textbf{Failure-signature clustering.} Failing executions are normalized ---
memory addresses, line numbers, temporary paths, timestamps, and UUIDs are
stripped --- and fingerprinted on the templated exception type, message, and top
stack frames, so runs that differ only in incidental detail collapse into a
single \emph{signature}. This is the unit on which root cause and risk are
assessed.

\textbf{Root cause and flaky-vs-real.} Each signature is categorized (null
dereference, timeout, dependency, configuration, resource, concurrency, network,
assertion, \ldots) with a confidence and the matched evidence. A signature is
flagged \emph{flaky} when a single commit exhibits both a pass and a failure
(order-agnostic flakiness) or a high pass/fail flip rate across runs; an
all-flaky failure set grades \emph{warn} rather than \emph{fail}, so flakiness
does not block the gate. A per-test risk score (recency-weighted failure rate
blended with instability) and a next-run failure forecast provide the detection
probability $q_t$ that the selection bound consumes.

\section{Selection with a Confidence Bound}
\label{sec:selection}

Let each candidate test $t$ carry a detection probability $q_t \in [0,1]$ --- the
modelled chance it fails (and thus catches a regression) on this run. If we run
subset $S$ and skip $U = \text{All} \setminus S$, then \emph{under an
independence assumption} the probability that at least one skipped test would
have failed is
\begin{equation}
\eps \;=\; 1 - \prod_{t \in U}\bigl(1 - q_t\bigr).
\label{eq:eps}
\end{equation}
We call $\eps$ the \emph{confidence lost} by skipping $U$: it is an upper-bound
proxy for the chance of missing a regression the suite would have caught.
Selection keeps the highest value-density tests ($q_t$ per unit time) until
either $\eps \le \eps_{\text{target}}$ or a time budget is exhausted; tests
touched by a code change can be force-kept. The achieved $\eps$ is always
reported and rounded \emph{up}, so the bound is never understated.

Naively evaluating Eq.~\eqref{eq:eps} at each greedy step is cubic in the number
of distinct tests; assaylab maintains a running product with a separate
zero-factor count (for $q_t = 1$), making selection $O(n \log n)$ after the sort
--- important because the candidate set is derived from an untrusted corpus and
must not become a compute-exhaustion vector.

\section{Attested, Re-derivable Receipts}
\label{sec:attest}

A selection emits a \emph{receipt} that commits, in a canonical serialization,
to: a hash of the candidate set and its $q$/duration inputs; the hashes and
sizes of the selected and skipped sets; the achieved $\eps$ and confidence; a
per-receipt nonce; and the tool version. The receipt is signed over exactly
these bytes, so the signature binds \emph{what actually happened}, not merely
that a run occurred. Two signing modes are provided:

\begin{itemize}[leftmargin=1.4em,itemsep=2pt]
  \item \textbf{HMAC-SHA256} (default) --- a symmetric trust domain: the verifier
        holds the key. The key is resolved from the environment or a persisted
        per-installation random key; there is no shipped default secret, and
        verification is constant-time.
  \item \textbf{ed25519} (asymmetric) --- the producer signs with a private key
        and publishes the public key; any external consumer verifies against the
        trusted public key \emph{without holding any secret}. This makes ``an
        external consumer verifies'' real rather than aspirational.
\end{itemize}

Crucially, a signature proves only that \emph{some key-holder committed} to an
$\eps$. Because selection is deterministic in the inputs the receipt commits to,
a consumer that holds the candidate inputs can \emph{recompute} the selection
and confirm the receipt reproduces --- the bound is genuine, not merely signed.
A receipt whose $\eps$ has been forged fails both the signature check and the
reproduction check.

\section{Security}
\label{sec:security}

The signing, key-handling, and acceptance-gating machinery has attack surface,
so it was subjected to an adversarial review: \textbf{five independent rounds},
each a fresh attempt to forge a receipt, weaken a bound, or exhaust a resource,
with fixes applied between rounds and the process continued until a full round
produced no new exploitable finding. Confirmed issues found and closed (each
pinned by a regression test that runs the exploit and shows it blocked) include:
a signing-key symlink/TOCTOU vector enabling forgery; an $O(n^3)$
compute-exhaustion path in selection reachable from an untrusted corpus (fixed to
$O(n\log n)$); a non-finite-float path that split the signed bytes from the
stored bytes; and, in the LLM-assisted module (below), a trust-boundary defect
where a hand-crafted proposal could weaken its own acceptance criteria --- closed
by signing proposals at generation and refusing to grade a tampered one. The
core HMAC construction, canonicalization, and constant-time verification survived
every round.

We do not claim the system is unexploitable. What is verified is narrower and
honest: every \emph{discovered} vector is closed and regression-pinned, and the
machinery survived five adversarial rounds ending empty. Residual risk that no
audit eliminates is stated in Section~\ref{sec:limits}.

\paragraph{LLM-assisted generation, gated.} assaylab can propose a regression
test from a signature, or a mitigation for a flaky one, using an LLM. It
\emph{never executes or applies the generated code}: a proposal is a dry-run
artifact, run only by a human or CI in their own sandbox, and acceptance flows
back through the verdict layer --- a generated test counts only if a real run
shows it reproduces the failure with the original signature; a heal counts only
if the flaky signature stops failing across distinct reruns.

\section{Evaluation}
\label{sec:eval}

We evaluate on the FlakeFlagger dataset~\cite{flakeflagger} (Zenodo 4450723, CC
BY 4.0): 26{,}765 tests across 23 Java projects, each rerun 10{,}000 times.

\subsection{The confidence bound holds on real data}

For each test we take its \emph{measured} per-run failure probability
(failing reruns over total reruns) as $q_t$ and run selection at a target
confidence loss of $\eps_{\text{target}} = 0.05$. Table~\ref{tab:bound} reports
the outcome. The selected subset yields a \textbf{35.9$\times$} reduction in
test-time, and the realized miss rate computed on the measured rates
($0.049957$) is \emph{within} the bound the receipt committed to ($0.049958$):
the receipt's $\eps$ is not a marketing figure but a statement that matches
reality on real failure data.

\begin{table}[h]
\centering
\begin{tabular}{lr}
\toprule
Tests & 26{,}765 \\
Target $\eps$ & 0.05 \\
Claimed $\eps$ (signed) & 0.049958 \\
Realized miss rate & 0.049957 \\
Test-time speedup & 35.9$\times$ \\
Bound holds (realized $\le$ claimed) & yes \\
\bottomrule
\end{tabular}
\caption{Confidence-bound validation on FlakeFlagger measured failure rates.}
\label{tab:bound}
\end{table}

\subsection{Flaky prediction from features (baseline)}

As a secondary check we train assaylab's built-in, dependency-free logistic
model on FlakeFlagger's feature table (70/30 held-out split; decision threshold
selected on the training split only, since the set is 0.7\% positive). It
attains precision $0.34$, recall $0.26$, $F_1 = 0.29$. This is a deliberately
lightweight baseline: FlakeFlagger's own tuned random forest reports higher
$F_1$ (${\approx}0.6$); assaylab trades accuracy for an inspectable,
JSON-serializable model that carries no heavy dependency and cannot execute code
on load. The point of the paper is the confidence bound above, for which the
flaky/real signal is an input.

\section{Limitations and Residual Risk}
\label{sec:limits}

The bound of Eq.~\eqref{eq:eps} rests on stated assumptions, not guarantees:
\emph{independence} of test failures; \emph{stationarity} of $q_t$ (history
predicts the next run); and coverage only of \emph{regression classes seen
historically}. Correlated failures or a novel regression class are outside the
model. On the attestation side: HMAC is a symmetric trust domain (mitigated by
the ed25519 mode); receipts are stateless, so freshness/replay must be enforced
consumer-side (a nonce ledger plus a max-age window); and the acceptance gate
trusts the run output a user supplies from their own sandbox. These are design
boundaries, honestly disclosed, not defects.

\section{Conclusion}

assaylab reframes test-suite reduction as an \emph{attestable} operation: it does
not merely run fewer tests, it emits a signed, independently re-derivable
statement of the confidence lost. On real measured failure data the bound holds
at a 35.9$\times$ speedup, and the machinery that produces it survived five
rounds of adversarial review. We hope the ``speedup with a checkable bound''
framing is useful beyond this implementation --- any selection tool can, in
principle, attach and sign a confidence statement rather than asking for trust.

\section*{Availability}

Source: \url{https://github.com/amitpatole/assaylab} (MIT). Package:
\url{https://pypi.org/project/assaylab/} (\code{pip install assaylab}).
Documentation: \url{https://amitpatole.github.io/assaylab/}. The software release
archived alongside this paper carries its own DOI, cross-referenced in the Zenodo
record.

\begin{thebibliography}{9}
\bibitem{idflakies} W. Lam, R. Oei, A. Shi, D. Marinov.
  \emph{iDFlakies: A Framework for Detecting and Partially Classifying Flaky
  Tests}. ICST, 2019.
\bibitem{flakeflagger} A. Alshammari, C. Morris, M. Hilton, J. Bell.
  \emph{FlakeFlagger: Predicting Flakiness Without Rerunning Tests}. ICSE, 2021.
  Dataset: Zenodo 4450723 (CC BY 4.0).
\bibitem{deflaker} J. Bell, O. Legunsen, M. Hilton, L. Eloussi, T. Yung,
  D. Marinov. \emph{DeFlaker: Automatically Detecting Flaky Tests}. ICSE, 2018.
\bibitem{rtptorrent} T. Mattis, P. Rein, F. D\"ursch, R. Hirschfeld.
  \emph{RTPTorrent: An Open-source Dataset for Evaluating Regression Test
  Prioritization}. MSR, 2020.
\bibitem{travistorrent} M. Beller, G. Gousios, A. Zaidman.
  \emph{TravisTorrent: Synthesizing Travis CI and GitHub for Full-Stack Research
  on Continuous Integration}. MSR, 2017.
\end{thebibliography}

\end{document}
